\documentclass[11pt,a4paper]{article}
\usepackage[utf8]{inputenc}
\usepackage[margin=2.4cm]{geometry}
\usepackage{amsmath,amssymb}
\usepackage{booktabs}
\usepackage{xcolor}
\usepackage{graphicx}
\usepackage{listings}
\usepackage{enumitem}
\usepackage{microtype}
\usepackage[colorlinks=true,linkcolor=blue!50!black,urlcolor=blue!50!black]{hyperref}

\definecolor{codegray}{gray}{0.45}
\definecolor{codeblue}{rgb}{0.1,0.2,0.6}
\lstset{
  language=Python,
  basicstyle=\ttfamily\footnotesize,
  keywordstyle=\color{codeblue},
  commentstyle=\color{codegray}\itshape,
  stringstyle=\color{teal},
  showstringspaces=false,
  breaklines=true,
  frame=single,
  rulecolor=\color{black!25},
  numbers=left,
  numberstyle=\tiny\color{codegray},
  literate={θ}{$\theta$}1 {ξ}{$\xi$}1,
}

\newcommand{\session}[3]{%
  \title{\textbf{#1}\\[2pt]\large #2}
  \author{astro-lab \,--\, #3}
  \date{}
  \maketitle
  \vspace{-1.5em}
}
\newcommand{\todo}[1]{\textcolor{red!60!black}{\textbf{[#1]}}}
\usepackage{framed}
\newenvironment{keyidea}{%
  \def\FrameCommand{\color{blue!40!black}\vrule width 2.5pt\hspace{8pt}\color{black}}%
  \MakeFramed{\advance\hsize-\width\FrameRestore}}{\endMakeFramed}

\newcommand{\dd}{\mathrm{d}}
\newcommand{\pc}{\,\mathrm{pc}}
\newcommand{\Msun}{M_\odot}
\newcommand{\Rsun}{R_\odot}
\setlist[itemize]{itemsep=2pt,topsep=4pt}
\setlist[enumerate]{itemsep=3pt,topsep=4pt}
\setlength{\parindent}{0pt}
\setlength{\parskip}{6pt}

\begin{document}
\session{Session 5: Exoplanets from Light Curves}{Transit geometry, detrending, the box-least-squares periodogram}{reading notes}

\textbf{Status:} reading notes and exercise specification. The code for this session (\texttt{05\_lightcurves/}) is not in the repository yet; the exercises below define what to build.

\section*{Goals}
Predict the depth and duration of a transit from physical parameters, then recover a planet's period from a noisy light curve with a box-least-squares search, first on synthetic data and then on real Kepler or TESS data via Lightkurve.

\section{Transit geometry}
A planet of radius $R_p$ crossing a star of radius $R_\star$ blocks a fraction of light
\[
\delta \approx \left(\frac{R_p}{R_\star}\right)^2 = k^2 .
\]
Examples: a Jupiter-sized planet crossing a Sun-like star gives $\delta\approx1\%$; an Earth-sized planet gives $\delta\approx84$ ppm.

For a circular orbit with impact parameter $b$ the total duration (first to fourth contact) is approximately
\[
T_{14}\approx\frac{P}{\pi}\,\frac{R_\star}{a}\sqrt{(1+k)^2-b^2},
\]
valid for $a\gg R_\star$. The semi-major axis follows from Kepler's third law (Session 1), $a^3=GM_\star P^2/4\pi^2$. For Earth around the Sun this gives $T_{14}\approx 13$ hours for $b=0$.

\begin{keyidea}
\textbf{Detection is a signal-to-noise problem.} With per-point noise $\sigma$ and $N_{\mathrm{in}}$ in-transit points, $\mathrm{SNR}\approx\delta\sqrt{N_{\mathrm{in}}}/\sigma$. Small planets need many transits and long baselines.
\end{keyidea}

\section{Preparing a light curve}
\begin{enumerate}
\item Remove bad-quality cadences and normalise by the median flux.
\item \textbf{Flatten} slow stellar and instrumental trends with a smoothing filter whose window is several times the expected transit duration; a window that is too short will erase the transit itself.
\item Optionally clip upward outliers (flares), never downward ones.
\end{enumerate}

\section{Box-least-squares (BLS)}
BLS models the light curve as a constant level with a periodic rectangular dip, fitted over a grid of trial periods $P$, durations $\tau$ and epochs $t_0$. At each trial period the data are phase-folded; for each $(\tau,t_0)$ the in-box mean is compared with the out-of-box mean, and the statistic reported is the reduction in $\chi^2$ (or SNR) relative to a flat model. The best peak in the periodogram gives $P$, $\tau$, $t_0$ and a depth.

Practical points:
\begin{itemize}
\item The period grid must be fine enough that phase drift over the baseline $T$ is smaller than the duration: $\Delta P \lesssim P\,\tau/T$.
\item Peaks also appear at $P/2$, $2P$, $P/3$; check by folding at each.
\item A single deep dip does not give a period; you need at least two transits.
\end{itemize}

\section{Implementation plan}
\begin{enumerate}
\item Simulate: flat light curve, 2-minute cadence over 27 days (one TESS sector), inject a box transit with $P=3.7$ d, $\delta=0.5\%$, $\tau=2.5$ h, add Gaussian noise of 0.1\%.
\item Search with \texttt{astropy.timeseries.BoxLeastSquares}; confirm it recovers $P$ within the grid resolution.
\item Reduce the injected depth until recovery fails; plot recovery rate versus SNR.
\item Repeat on a real target with Lightkurve: download a light curve, flatten it, compute a BLS periodogram with \texttt{to\_periodogram}, and fold at the best period.
\end{enumerate}

\section{Exercises}
\begin{enumerate}
\item Verify the Earth and Jupiter depths and compute $T_{14}$ for both around the Sun.
\item Show that the transit probability for a random orientation is $\approx R_\star/a$ and evaluate it for $P=3.7$ d and $P=365$ d.
\item Inject the transit of a hot Jupiter into a real, quiet star's light curve and test whether your pipeline recovers it.
\item Compare BLS with a Lomb--Scargle periodogram on the same data and explain why Lomb--Scargle is poorly suited to box-shaped dips.
\end{enumerate}

\section*{References}
\begin{itemize}
\item Winn (2010), \textit{Transits and Occultations}, in \textit{Exoplanets} (arXiv:1001.2010)
\item Kov\'acs, Zucker, Mazeh (2002), \textit{A\&A} 391, 369 (BLS)
\item Lightkurve documentation: \url{https://lightkurve.github.io/lightkurve/}
\end{itemize}
\end{document}
